\documentclass{umk-matematika}

\begin{document}

\maketitlepage
    {Практическая работа №21}
    {Пирамида. Площадь поверхности и объём}
    {}

\section{Фундамент (1--4 балла)}

\textbf{Задача 1 (1 балл).} В правильной треугольной пирамиде сторона основания $6$ см, высота $10$ см. Найдите объём.

\textbf{Задача 2 (2 балла).} Объём пирамиды $48$ см³, площадь основания $12$ см². Найдите высоту.

\textbf{Задача 3 (3 балла).} В правильной четырёхугольной пирамиде сторона основания $4$ см, апофема $5$ см. Найдите площадь боковой поверхности.

\textbf{Задача 4 (4 балла).} Основание пирамиды — квадрат со стороной $6$ см. Высота пирамиды $9$ см. Найдите объём.

\section{Стены (5--7 баллов)}

\textbf{Задача 5 (5 баллов).} В правильной треугольной пирамиде сторона основания $4$ см. Боковое ребро наклонено к плоскости основания под углом $60^\circ$. Найдите объём.

\textbf{Задача 6 (6 баллов).} В правильной четырёхугольной пирамиде высота $8$ см, сторона основания $12$ см. Найдите площадь полной поверхности.

\textbf{Задача 7 (7 баллов).} Основание пирамиды — ромб с диагоналями $6$ см и $8$ см. Высота пирамиды $5$ см. Найдите объём.

\section{Кровля (8--10 баллов)}

\textbf{Задача 8 (8 баллов).} Крыша башни имеет форму правильной четырёхугольной пирамиды. Сторона основания $8$ м, высота $3$ м. Сколько листов кровельного материала потребуется, если один лист покрывает $2$ м²? (Найдите площадь боковой поверхности с запасом $10\%$).

\textbf{Задача 9 (9 баллов).} Пирамида имеет форму правильной четырёхугольной пирамиды со стороной основания $300$ м и высотой $150$ м. Найдите её объём в м³.

\textbf{Задача 10 (10 баллов).} Памятник состоит из куба со стороной $2$ м и правильной четырёхугольной пирамиды, поставленной сверху. Сторона основания пирамиды $2$ м, высота $1{,}5$ м. Найдите общий объём памятника.

\newpage
\section*{Ответы}

\begin{enumerate}
  \item $30\sqrt{3}$ см³
  \item $12$ см
  \item $40$ см²
  \item $108$ см³
  \item $\frac{16\sqrt{3}}{3}$ см³
  \item $384$ см²
  \item $40$ см³
  \item $44$ листа
  \item $4500000$ м³ ($4{,}5$ млн м³)
  \item $10$ м³
\end{enumerate}

\end{document}